\documentclass[11pt,a4paper]{article}

\usepackage[T1]{fontenc}
\usepackage[utf8]{inputenc}
\usepackage{lmodern}
\usepackage[margin=2.5cm]{geometry}
\usepackage{microtype}
\usepackage{enumitem}
\usepackage{xcolor}
\usepackage{hyperref}

\setlength{\parindent}{0pt}
\setlength{\parskip}{0.6em}
\setlist[itemize]{topsep=0.2em,itemsep=0.25em}

\definecolor{linkblue}{HTML}{1F4E79}
\hypersetup{
    colorlinks=true,
    linkcolor=linkblue,
    urlcolor=linkblue,
    pdftitle={VIP Week 10 Report},
    pdfauthor={Yorgo Azzi}
}

\begin{document}

\begin{center}
    {\Large\bfseries VIP Week 10 Report}\\[0.5em]
    {\large Yorgo Azzi}\\[0.25em]
    November 24, 2026
\end{center}

\section{Overview}

This week, I refined the movement-controlled Conductor mode to give the
performer more control over the soprano melody and the sound of the three
generated accompaniment voices. The right hand still commands the beat and
selects chords, while the left hand performs the melody. I also added a button
that exchanges these roles during a performance.

\section{Work Completed}

\subsection{More Expressive Melody Gestures}

Vertical movement of the melody hand still determines whether the soprano
ascends or descends. Finger gestures now determine how far each movement goes:
the index finger moves one scale note, index and middle move two, and adding
more fingers increases the jump to three, four, or five scale notes. This lets
the performer move between nearby notes and larger leaps without changing the
underlying Markov accompaniment.

The thumb alone now plays an arpeggio through the current chord. For example,
starting on C over a C-major chord, upward movements produce E, G, and the
next C; downward movements produce the G, E, and C below. Tilting the melody
hand right or left plays one semitone up or down from the current pitch. The
hand must return near upright before another tilt step is accepted, reducing
repeated notes caused by tracking noise.

A closed fist now mutes only the soprano. This allows the performer to return
the hand to its starting height without sounding descending notes. Opening
the fist restarts the note from which the arpeggio began, ready for another
run. The earlier horizontal-position register control was removed to keep the
gesture mapping simpler.

\subsection{Hand Distance and Role Switching}

The distance between the hands now controls two aspects of the output.
Moving the hands apart increases MIDI expression, so even a held note can
grow louder, and encourages wider spacing between the accompaniment voices
at the next generated beat. Bringing the hands together produces a quieter,
more compact sound.

A \textbf{Swap hand roles} button and keyboard shortcut exchange the musical
functions of the hands. The new beat hand controls the conducting stroke and
chord fingers, while the other hand controls the melody. Motion tracking is
reset when the roles change so the switch itself does not create a false beat.
The same camera and MediaPipe tracker continue running.

\section{Results}

\begin{itemize}
    \item The performer can choose between short scale movements, larger
    melodic jumps, chord arpeggios, and semitone adjustments.
    \item The fist provides a silent way to reposition the melody hand before
    repeating an arpeggio.
    \item Hand distance affects held-note loudness immediately and influences
    accompaniment spacing when the next chord is generated.
    \item The performer can switch which hand controls the beat and which
    controls the melody without restarting the application.
    \item Automated gesture, audio-event, interface, and music-generation
    checks passed. The new mappings still need evaluation with a physical
    webcam and performer.
\end{itemize}

\section{Challenges}

Several gestures depend on reliable MediaPipe finger and palm-angle estimates.
In particular, the thumb must be distinguished from a closed fist, and a
deliberate tilt must be distinguished from small changes in hand posture.
The fist gesture also needs to mute promptly enough that returning the hand
to its starting position feels natural. These details require live testing
and adjustment of thresholds.

Hand distance currently controls both loudness and accompaniment spread.
Although both changes suggest a larger musical sound, performer trials are
needed to determine whether this combined mapping is clear and useful.

\section{Next Steps}

\begin{itemize}
    \item Explore gestures in which both hands move together, and test whether
    the distance between them can determine whether the movement controls a
    chord or the melody.
    \item Profile the visual rendering path and investigate faster drawing,
    including OpenGL or Vulkan where appropriate, so visual feedback keeps up
    with the performance.
    \item Add 3D particle effects that respond to musical events and hand
    movement, while keeping the camera and audio response responsive.
\end{itemize}

\end{document}
